\newcommand{\subrec}{RECLAMAÇÃO DE COBRANÇA INDEVIDA SOBRE ENQUADRAMENTO TARIFÁRIO INCORRETO}
\section{DOS FATOS E DOS DIREITOS}

\begin{enumerate}[resume]
\item
Verificou-se que a UC \unidadeConsumidora{} está cadastrada com enquadramento tarifário errado, gerando, desta forma, prejuízos ao consumidor, por erro de cadastro por parte da concessionária.
\item  
Assim, a UC \unidadeConsumidora{} deve ser cadastrada conforme art. 191, inciso I, da Resolução Normativa 1000/2021 da Aneel, vale dizer, deve ser enquadrada na classe serviço público, subclasse água, esgoto e saneamento.  
\item 
A seguir transcreve-se o art. 191 da citada resolução:

\Citacao{Art. 191. Deve ser classificada na classe serviço público a unidade consumidora de responsabilidade do poder público ou daquele que receba essa delegação, destinada exclusivamente ao fornecimento de energia elétrica para motores, máquinas e cargas essenciais à operação de serviços públicos nas seguintes subclasses: 

I - água, esgoto e saneamento; e

II - tração elétrica.}

\item 
Vale ressaltar que essa classificação já era definida desde a REN 414/2010, que foi revogada pela REN 1000/2021, conforme trecho a seguir:

\Citacao{Art. 53-Q Na classe serviço público enquadram-se as unidades consumidoras que se destinem, exclusivamente, ao fornecimento para motores, máquinas e cargas essenciais à operação de serviços públicos de água, esgoto, saneamento e tração elétrica urbana ou ferroviária, explorados diretamente pelo Poder Público ou mediante concessão ou autorização, considerando-se as seguintes subclasses:

I – tração elétrica; e 

II – água, esgoto e saneamento.}

\item 
A UC relacionada acima tem direito a benefício tarifário, conforme art. 664 da citada 
resolução, transcrito a seguir:

\Citacao{Art. 664. A unidade consumidora classificada na subclasse água, esgoto e saneamento, conforme disposições do Decreto Nº 7.891, de 2013, tem direito ao benefício de redução nas tarifas aplicáveis, nos percentuais a seguir:

I - 2021: redução de 6\% (seis por cento);

II - 2022: redução de 3\% (três por cento); e

III - 2023: sem redução.}

\item  
Vale destacar que inicialmente o desconto aplicado era de 15\%, conforme REN 414/2010:

\Citacao{Art. 53-R. As unidades consumidoras classificadas na subclasse água, esgoto e 
saneamento, conforme disposições do Decreto Nº 7.891, de 2013, tem direito ao 
benefício tarifário de redução nas tarifas aplicáveis, nos percentuais da tabela a 
seguir: (Incluído pela REN ANEEL 800, de 19.12.2017)

\vspace{0.4cm}

\renewcommand{\arraystretch}{1.8} % Controlar altura das linhas
\begin{tabular}{|c|c|c|c|c|}
\hline
Grupo & TUSD R\$/kW & TUSD R\$/MWh & TE R\$/MWh & Tarifa para aplicação da redução \\
\hline 
A & 15\% & 15\% & 15\% & tarifas azul e verde \\
\hline 
B & -- & 15\% & 15\% & B3\\
\hline
\end{tabular}
}

\item 
Contudo, conforme decreto Nº 9.642 de 27 de dezembro de 2018 houve previsão da redução gradual dos descontos tarifários. Dessa forma, até o final de 2018 o desconto era fixo em 15\%.

\Citacao{§ 4º A partir de 1º de janeiro de 2019, nos respectivos reajustes ou procedimentos ordinários de revisão tarifária, os descontos de que trata o § 2º serão reduzidos à razão de vinte por cento ao ano sobre o valor inicial, até que a alíquota seja zero.” (NR)}

\end{enumerate}

\section{DOS PEDIDOS}

\begin{enumerate}[resume]
\item 
Diante do exposto, requer que a Concessionária:

\begin{enumerate}
\item 
Efetue vistoria em tal UC e proceda com a correção do enquadramento tarifário da mesma em conformidade com o estabelecido na Resolução Normativa Nº 1.000/2021 da ANEEL.
\item 
Efetue a devolução de todos os valores cobrados a maior, referente aos últimos 10 (dez) anos. Conforme o Despacho Nº 2.006 da ANEEL de 08 de julho de 2024, observando o prazo prescricional de 10 anos, conforme previsto no art. 205 do Código Civil.
\item 
Devolva os valores cobrados a maior em dobro, atualizados monetariamente pelo IPCA e juros de mora de 1\% ao mês calculados pro \textit{rata die}, conformidade com o estabelecido na Resolução Normativa Nº 1.000/2021 da ANEEL.   
\item 
Que apresente as seguintes informações: Data de ligação da UC, se houve modificação de tarifa entre a data da ligação e a data da resposta.   
\item 
Que em caso de indeferimento apresente os argumentos e as provas (comprovações) para não enquadramento (fotos, carga instalada etc).   
\item 
Proceda a análise e atendimento a este requerimento no prazo máximo de 10 (dez) dias úteis, para informar o resultado da classificação tarifária no caso de necessidade de vistoria, em conformidade com o estabelecido na Resolução Normativa Nº 1.000/2021 da ANEEL.  
\item 
Proceda com a devolução do valor integral, por meio de depósito bancário, na conta corrente do Município, conforme estabelece o § 7º, do art. 323 da REN Nº 1.000/2021 da ANEEL, comunicando por escrito o resultado do deferimento desta reclamação, por meio de e-mail.

\end{enumerate}
\end{enumerate}